\scanpage{187}
\mspage{96}
% L01–03
peut aussi \deleted{formuler} \inserted{expliciter} en termes de théorie des groupes. Notons que l'on aura alors
% slanted
\margin{il faudra y revenir par la suite\ldots}
\[
\operatorname{Et}_{A_0}(s)=\operatorname{Et}_{A_\lambda}(\lambda s)\tag*{(26)}
\]

% L03 margin
% slanted
\margin{Cette relation est indépendante de la condition que l'homothétie soit stricte -- mais si elle est stricte, $A_\lambda$ étant associée à une structure de polyèdre, la combinatoire induite par $A_\lambda$ sur une étoile $\inserted{E}$ de $A_\lambda$, i.e. une structure $A_{E,\lambda}\subset\unsure{\mathfrak P_2(E)|\mathrm P_2(E)}$, est donc une structure de polygône -- ce qui ne semble peut être pas vrai autrement\ldots}

% L04–08
$s\in S$, donc les deux structures de graphe $A_0$ et $A_\lambda$ donnent lieu aux mêmes étoiles -- mais avec des centres différents : Le centre $\lambda s$ de l'étoile $E$, pour la structure $\Gamma_\lambda=\underset{2}{(}S,A_\lambda)$, est l'antipodique du centre $s$ pour $\Gamma_0$, via l'antipodisme $\lambda$.

% L09–19
Considérons alors le groupe $G_{0\,s}=G_{0\,(s'=\lambda.s)}$, \deleted{et le sous} le sous-groupe des rotations relatif à \deleted{\unsure{leur}} $s$ est le seul sous-groupe cyclique d'ordre $n$ de $G_{0\,s}\simeq\mathbb D_n$, donc est égal au sous-groupe correspondant \deleted{pour} relatif à $s'$, Cela montre que sur l'étoile $s$, les structures de polygône, induites par $A_0$ et $A_\lambda$ respectivement, ont $=$ groupe de rotations, i.e. sont transformées l'une de l'autre par un élément \inserted{$\gamma$} de $(\mathbb Z/n\mathbb Z)^*/\{\pm1\}=Z'_n$ ($n$ étant l'ordre des points de $S$) :
\[
\alpha_{E,\lambda}=(\alpha_{E,0})^{(\gamma)},\qquad\gamma\in Z'_n=(\mathbb Z/n\mathbb Z)^*/\pm1\tag*{(27)}
\]

% archive sheet number: 186
